\documentclass{article}
\usepackage{amsmath}
\usepackage{amsthm}
\usepackage{amsfonts}
\usepackage{amssymb}
\usepackage{bbold}
\usepackage{dsfont}
\usepackage{graphicx}
\usepackage{cleveref}
\usepackage{bm}
\usepackage{booktabs}
\usepackage{siunitx}
\usepackage{xcolor}
\usepackage{pdfpages}
\usepackage[utf8]{inputenc}
\usepackage[english]{babel}
\newtheorem{definition}{Definition}
\newtheorem{example}{Example}
\newcommand{\exP}[1]{\mathrm{e}^{#1}}
\newcommand{\dif}{\mathrm{d}}
\renewcommand{\vec}[1]{\mathbf{#1}}
\newcommand{\pdagger}{{\phantom{\dagger}}}
\newcommand{\state}[1]{\vert #1\rangle}
\newcommand{\cstate}[1]{\langle #1\vert}
\newcommand{\tr}[1]{\mathrm{tr}\left\{#1\right\}}
\newcommand{\trs}[2]{\, \mathrm{tr}_{#1}\{#2\}}
\newcommand{\expc}[3]{\langle #1\vert #2 \vert #3\rangle}
\usepackage[doi=false,url=false,isbn=false,backend=bibtex,style=nature,natbib=true,sorting=none]{biblatex}
\bibliography{bib.bib}
\usepackage{geometry}
\geometry{
	a4paper,
	total={170mm,257mm},
	left=15mm,
	right=15mm,
	top=15mm,
}
\begin{document}
\title{Mapping approach to dynamics in the spinless Peierls model}
\maketitle
\section{Theory}
The Peierls model consists of electronic sites and bond vibrations, which couple to electrons via nearest neighbor hoppings
\begin{equation}
	H = \varepsilon_0 \sum_j  d^\dagger_j d_j^\pdagger + \omega_0 \sum_j a_j^\dagger a_j^\pdagger + \sum_{j}' (d_j^\dagger d_{j+1} + \mathrm{h.c.})( -t_0 + t_1(a_j^\dagger + a_j)).
\end{equation}
The first approximation is the semi-classical limit for the vibrations, which can be achieved, for example, via the correspondence principle. The resulting model is then
\begin{equation}
	H/{\hbar\omega_0} = \varepsilon_0 \sum_j  d^\dagger_j d_j^\pdagger + \frac{1}{2} \sum_j (x_j^2+p_j^2) + \sum_{j}' (d_j^\dagger d_{j+1} + \mathrm{h.c.})( -t_0 + \frac{t_1}{\sqrt{2}}x_j).
\end{equation}
where the parameters $\epsilon_0,t_0,t_1$ are all dimensionless.
Next, we employ a mapping of the fermionic operators
\begin{align}
	c_j\mapsto \frac{1}{\sqrt{2}} (X_j+iP_j)\\
	c_j^\dagger \mapsto \frac{1}{\sqrt{2}} (X_j-iP_j)
\end{align}
which results in the final Hamiltonian
\begin{equation}
	H = \frac{\epsilon_0}{2} \sum_j  (X_j^2+P_j^2-2\gamma) + \frac{1}{2} \sum_j (x_j^2+p_j^2) + \frac{1}{2}\sum_{j}' (X_j X_{j+1} + P_j P_{j+1})( -t_0 + \frac{t_1}{\sqrt{2}}x_j).
\end{equation}
where $\gamma=\frac{\sqrt{N+1}-1}{N}$ is the zero-point energy parameter as a function of the number of electronic sites $N$. 
The equations of motion generated by this Hamiltonian are
\begin{align}
	\dot{x}_j &= p_j \\
	\dot{p}_j &= -x_j - \frac{t_1}{2\sqrt{2}}(X_j X_{j+1} + P_j P_{j+1}) \\
	\dot{X}_j &= \varepsilon_0 P_j + \frac{P_{j+1}}{2}(-t_0 + \frac{t_1}{\sqrt{2}}x_j)\\
	\dot{P}_j &= -\varepsilon_0 X_j - \frac{X_{j+1}}{2} (-t_0+\frac{t_1}{\sqrt{2}}x_j).
\end{align}
At initial time, the system is factorized with a single electronic excitation on site $1$ and all other sites unoccupied, which is realized by sampling on points on the circle in the mapping space of each site
\begin{align*}
	n_j = \langle c_j^\dagger c_j^\pdagger\rangle = \frac{1}{2}(X_j^2+P_j^2 -2\gamma) =\delta_{j1}.
\end{align*}
The vibrations are assumed to be thermal, sampled from the Wigner distribution 
\begin{equation}
	\rho_\mathrm{vib}(x,p) = \prod_j 2\alpha_j \exP{-\alpha_j (x_j^2+p_j^2)}
\end{equation} 
with $\alpha_j = \tanh(\frac{\beta \hbar\omega_j}{2})$. 
These equations of motion are solved using the momentum integral (MInt) algorithm \cite{cook2023algorithm}. Therein, one splits the Hamiltonian into $H_1$ and $H_2$
\begin{align}
	H_1 &= \frac{1}{2}\sum_j p_j^2\\
	H_2 &= \frac{1}{2}(P^T V(x)P + X^T V(x) X - \tr{V(x)}).
\end{align}
The algorithm 
\clearpage


\end{document}


